\section{Conclusion}
OpenVisualJev combines tree-masked multimodal encoding, position-free set interaction, dynamic pointer scoring, calibration-aware training, and held-out temperature scaling into a multimodal decision pipeline. The design combines permutation-equivariant scores with preferences conditioned on the candidate set. Exhaustive permutations verify decision stability, and preference-reversal tasks demonstrate the expressive role of candidate interaction. Under the evaluated single-seed balanced protocol, Set achieves higher macro accuracy than the complete-candidate LM across six recognition datasets; dynamic candidate-count training improves both models on larger candidate sets. Separately supervised natural VQA yields accuracy close to the native LM. The balanced training-loss ablation supports combining mean-confidence regularization with temperature scaling to improve probability quality. These findings identify candidate-set structure as a useful design principle for multimodal decision models and provide an empirical basis for its use in dynamic choice tasks.
